% GENERATED FILE. Edit canonical lessons, then run npm run generate:books.
% canonical-source-hash: fnv1a64:f61a70291fffbae8
% canonical-lessons: RU-C334-farsh, RU-C334-farshirovat, RU-C334-protsezhivat, RU-C334-durshlag, RU-C334-sito, RU-R334-review-327-334

\chapter{Stuffing and straining}
\label{ch:ru-stuffing-and-straining}
\hlchaptermodality{334}

\begin{chapteropening}
\textbf{By the end of this chapter:} \emph{I can talk about mince and stuffing, and about straining with a colander and a sieve.}

Complete the last lesson of chapter 334: I can talk about mince and stuffing, and about straining with a colander and a sieve.
\end{chapteropening}

\section[\texorpdfstring{\ru{фарш} (farsh)}{фарш (farsh)}]{\ru{фарш} (farsh) — mince, minced meat}

\label{lesson:RU-C334-farsh}



\begin{quote}
Before the new one: say the Russian for a seasoning, then the Russian for to marinate.
\end{quote}

\subsection*{You'll want to know: \ru{фарш}}
\textbf{\ru{фарш}} (\emph{farsh}) — \textquotedblleft{}mince, minced meat\textquotedblright{}.

\begin{quote}
\ru{Для} \ru{котлет} \ru{нужен} \ru{фарш}.
\end{quote}

\begin{quote}
\emph{Dlya kotlét núzhen farsh.}
\end{quote}

\begin{quote}
You need mince for meat patties.
\end{quote}

Masculine.

\begin{cousinweb}
\emph{\ru{фарш}} was borrowed from French \emph{farce} \textquotedblleft{}stuffing\textquotedblright{} (EtymDB), from Latin \emph{farcīre} \textquotedblleft{}to cram, to stuff\textquotedblright{}.
\end{cousinweb}

\begin{grammarlens}[title={What you've built}]
One more word for telling, explaining and describing at B1.
\end{grammarlens}

\subsection*{Your turn}
\begin{itemize}
\raggedright
  \item \emph{Say it:} \emph{farsh}
  \item \emph{Say it:} \emph{farsh}, once more
  \item \emph{Say it:} a sentence of your own with \emph{farsh}
  \item \emph{Recall:} say the Russian for a seasoning, then the Russian for to marinate, then say \emph{farsh} again
\end{itemize}

\begin{tcolorbox}[breakable,colback=teal!4,colframe=teal!35!black,title={Before you move on}]
What does \ru{фарш} mean? (\textquotedblleft{}Mince, minced meat\textquotedblright{}.) Say it once more.
\end{tcolorbox}
\section[\texorpdfstring{\ru{фаршировать} (farshirovát)}{фаршировать (farshirovát)}]{\ru{фаршировать} (farshirovát) — to stuff}

\label{lesson:RU-C334-farshirovat}



\begin{quote}
Before the new one: say the Russian for to marinate, then the Russian for mince.
\end{quote}

\subsection*{You'll want to know: \ru{фаршировать}}
\textbf{\ru{фаршировать}} (\emph{farshirovát}) — \textquotedblleft{}to stuff\textquotedblright{}.

\begin{quote}
\ru{Мама} \ru{фарширует} \ru{перец} \ru{рисом} \ru{и} \ru{мясом}.
\end{quote}

\begin{quote}
\emph{Máma farshirúyet pérets rísom i myásom.}
\end{quote}

\begin{quote}
Mum stuffs peppers with rice and meat.
\end{quote}

Imperfective; perfective \emph{\ru{нафаршировать}} (\emph{nafarshirovát}). The filling is in the instrumental.

\begin{cousinweb}
\emph{\ru{фаршировать}} is built on \emph{\ru{фарш}} \textquotedblleft{}stuffing\textquotedblright{}, from French \emph{farce}, Latin \emph{farcīre} \textquotedblleft{}to stuff\textquotedblright{} (EtymDB), with the verb ending \emph{-\ru{ировать}}.
\end{cousinweb}

\begin{grammarlens}[title={What you've built}]
One more word for telling, explaining and describing at B1.
\end{grammarlens}

\subsection*{Your turn}
\begin{itemize}
\raggedright
  \item \emph{Say it:} \emph{farshirovát}
  \item \emph{Say it:} \emph{farshirovát}, once more
  \item \emph{Say it:} a sentence of your own with \emph{farshirovát}
  \item \emph{Recall:} say the Russian for to marinate, then the Russian for mince, then say \emph{farshirovát} again
\end{itemize}

\begin{tcolorbox}[breakable,colback=teal!4,colframe=teal!35!black,title={Before you move on}]
What does \ru{фаршировать} mean? (\textquotedblleft{}To stuff\textquotedblright{}.) Say it once more.
\end{tcolorbox}
\section[\texorpdfstring{\ru{процеживать} (protsézhivat)}{процеживать (protsézhivat)}]{\ru{процеживать} (protsézhivat) — to strain}

\label{lesson:RU-C334-protsezhivat}



\begin{quote}
Before the new one: say the Russian for mince, then the Russian for to stuff.
\end{quote}

\subsection*{You'll want to know: \ru{процеживать}}
\textbf{\ru{процеживать}} (\emph{protsézhivat}) — \textquotedblleft{}to strain\textquotedblright{}.

\begin{quote}
\ru{Он} \ru{процедил} \ru{бульон}.
\end{quote}

\begin{quote}
\emph{On protsedíl bulón.}
\end{quote}

\begin{quote}
He strained the broth.
\end{quote}

Imperfective; perfective \emph{\ru{процедить}} (\emph{protsedít}), \emph{\ru{процежу}}.

\begin{cousinweb}
\emph{\ru{процеживать}} is built on \emph{\ru{про}-} \textquotedblleft{}through\textquotedblright{} and \emph{\ru{цедить}} \textquotedblleft{}to strain\textquotedblright{}, inherited from Proto-Slavic \emph{cěditi} \textquotedblleft{}to strain, to filter\textquotedblright{} (EtymDB).
\end{cousinweb}

\begin{grammarlens}[title={What you've built}]
One more word for telling, explaining and describing at B1.
\end{grammarlens}

\subsection*{Your turn}
\begin{itemize}
\raggedright
  \item \emph{Say it:} \emph{protsézhivat}
  \item \emph{Say it:} \emph{protsézhivat}, once more
  \item \emph{Say it:} a sentence of your own with \emph{protsézhivat}
  \item \emph{Recall:} say the Russian for mince, then the Russian for to stuff, then say \emph{protsézhivat} again
\end{itemize}

\begin{tcolorbox}[breakable,colback=teal!4,colframe=teal!35!black,title={Before you move on}]
What does \ru{процеживать} mean? (\textquotedblleft{}To strain\textquotedblright{}.) Say it once more.
\end{tcolorbox}
\section[\texorpdfstring{\ru{дуршлаг} (durshlág)}{дуршлаг (durshlág)}]{\ru{дуршлаг} (durshlág) — a colander}

\label{lesson:RU-C334-durshlag}



\begin{quote}
Before the new one: say the Russian for to stuff, then the Russian for to strain.
\end{quote}

\subsection*{You'll want to know: \ru{дуршлаг}}
\textbf{\ru{дуршлаг}} (\emph{durshlág}) — \textquotedblleft{}a colander\textquotedblright{}.

\begin{quote}
\ru{Макароны} \ru{откинули} \ru{на} \ru{дуршлаг}.
\end{quote}

\begin{quote}
\emph{Makaróny otkínuli na durshlág.}
\end{quote}

\begin{quote}
The pasta was drained in a colander.
\end{quote}

Masculine.

\begin{cousinweb}
\emph{\ru{дуршлаг}} was borrowed from German \emph{Durchschlag} (EtymDB), from \emph{durch} \textquotedblleft{}through\textquotedblright{} and \emph{Schlag} \textquotedblleft{}a blow\textquotedblright{}.
\end{cousinweb}

\begin{grammarlens}[title={What you've built}]
One more word for telling, explaining and describing at B1.
\end{grammarlens}

\subsection*{Your turn}
\begin{itemize}
\raggedright
  \item \emph{Say it:} \emph{durshlág}
  \item \emph{Say it:} \emph{durshlág}, once more
  \item \emph{Say it:} a sentence of your own with \emph{durshlág}
  \item \emph{Recall:} say the Russian for to stuff, then the Russian for to strain, then say \emph{durshlág} again
\end{itemize}

\begin{tcolorbox}[breakable,colback=teal!4,colframe=teal!35!black,title={Before you move on}]
What does \ru{дуршлаг} mean? (\textquotedblleft{}A colander\textquotedblright{}.) Say it once more.
\end{tcolorbox}
\section[\texorpdfstring{\ru{сито} (síto)}{сито (síto)}]{\ru{сито} (síto) — a sieve}

\label{lesson:RU-C334-sito}



\begin{quote}
Before the new one: say the Russian for to strain, then the Russian for a colander.
\end{quote}

\subsection*{You'll want to know: \ru{сито}}
\textbf{\ru{сито}} (\emph{síto}) — \textquotedblleft{}a sieve\textquotedblright{}.

\begin{quote}
\ru{Бабушка} \ru{просеяла} \ru{муку} \ru{через} \ru{сито}.
\end{quote}

\begin{quote}
\emph{Bábushka proséyala múku chérez síto.}
\end{quote}

\begin{quote}
Grandma sifted the flour through a sieve.
\end{quote}

Neuter.

\begin{cousinweb}
\emph{\ru{сито}} is inherited through Old East Slavic \emph{sito} from Proto-Slavic \emph{sito} \textquotedblleft{}sieve\textquotedblright{} (EtymDB).
\end{cousinweb}

\begin{grammarlens}[title={What you've built}]
One more word for telling, explaining and describing at B1.
\end{grammarlens}

\subsection*{Your turn}
\begin{itemize}
\raggedright
  \item \emph{Say it:} \emph{síto}
  \item \emph{Say it:} \emph{síto}, once more
  \item \emph{Say it:} a sentence of your own with \emph{síto}
  \item \emph{Recall:} say the Russian for to strain, then the Russian for a colander, then say \emph{síto} again
\end{itemize}

\begin{tcolorbox}[breakable,colback=teal!4,colframe=teal!35!black,title={Before you move on}]
What does \ru{сито} mean? (\textquotedblleft{}A sieve\textquotedblright{}.) Say it once more.
\end{tcolorbox}
\section[(dialogue)]{Words from chapters 327–334 — a review}

\label{lesson:RU-R334-review-327-334}



\begin{quote}
Say the Russian for an archaeologist and a dig.
\end{quote}

\subsection*{Your turn}
Say these aloud:

\begin{itemize}
\raggedright
  \item an archaeologist, a dig, a find, a treasure, a buried treasure
  \item a slice, a little piece, to slice, to grate, a grater
  \item to bring to the boil, to boil, to cook in water, a lid, to cover
  \item to mix, to whisk, a whisk, to stir, a ladle
  \item dough, to knead, to roll out, to shape, a filling
  \item to bake, a baking tray, to heat up, to get burnt, to brown
  \item to salt, to pepper, to sprinkle, a seasoning, to marinate
  \item mince, to stuff, to strain, a colander, a sieve
\end{itemize}

\begin{tcolorbox}[breakable,colback=teal!4,colframe=teal!35!black,title={Before you move on}]
An archaeologist? (\textbf{\ru{археолог}}.) A dig? (\textbf{\ru{раскопки}}.) A find? (\textbf{\ru{находка}}.) A treasure? (\textbf{\ru{сокровище}}.) A buried treasure? (\textbf{\ru{клад}}.) A slice? (\textbf{\ru{ломтик}}.) A little piece? (\textbf{\ru{кусочек}}.) To slice? (\textbf{\ru{нарезать}}.) To grate? (\textbf{\ru{натирать}}.) A grater? (\textbf{\ru{тёрка}}.) To bring to the boil? (\textbf{\ru{кипятить}}.) To boil? (\textbf{\ru{кипеть}}.) To cook in water? (\textbf{\ru{отваривать}}.) A lid? (\textbf{\ru{крышка}}.) To cover? (\textbf{\ru{накрывать}}.) To mix? (\textbf{\ru{смешивать}}.) To whisk? (\textbf{\ru{взбивать}}.) A whisk? (\textbf{\ru{венчик}}.) To stir? (\textbf{\ru{помешивать}}.) A ladle? (\textbf{\ru{половник}}.) Dough? (\textbf{\ru{тесто}}.) To knead? (\textbf{\ru{месить}}.) To roll out? (\textbf{\ru{раскатывать}}.) To shape? (\textbf{\ru{лепить}}.) A filling? (\textbf{\ru{начинка}}.) To bake? (\textbf{\ru{запекать}}.) A baking tray? (\textbf{\ru{противень}}.) To heat up? (\textbf{\ru{разогревать}}.) To get burnt? (\textbf{\ru{подгорать}}.) To brown? (\textbf{\ru{обжаривать}}.) To salt? (\textbf{\ru{солить}}.) To pepper? (\textbf{\ru{перчить}}.) To sprinkle? (\textbf{\ru{посыпать}}.) A seasoning? (\textbf{\ru{приправа}}.) To marinate? (\textbf{\ru{мариновать}}.) Mince? (\textbf{\ru{фарш}}.) To stuff? (\textbf{\ru{фаршировать}}.) To strain? (\textbf{\ru{процеживать}}.) A colander? (\textbf{\ru{дуршлаг}}.) A sieve? (\textbf{\ru{сито}}.)
\end{tcolorbox}
